\chapter*{Appendix C: Classification of Surgical Wounds}
\addcontentsline{toc}{chapter}{Appendix C: Classification of Surgical Wounds}
\markboth{Appendix C: Surgical Wounds}{Appendix C: Surgical Wounds}

The Centers for Disease Control and Prevention (CDC) classifies surgical wounds based on the degree of microbial contamination at the time of operation. This classification system predicts the risk of postoperative Surgical Site Infection (SSI) and guides the administration of prophylactic antibiotics.

\section*{Class I: Clean Wounds}
\textbf{Definition:} An uninfected operative wound in which no inflammation is encountered and the respiratory, alimentary, genital, or uninfected urinary tracts are not entered.
\begin{itemize}
  \item \textbf{Criteria:}
  \begin{itemize}
    \item Closed primarily (or closed drainage system).
    \item No break in aseptic technique.
    \item Non-traumatic, elective procedures.
  \end{itemize}
  \item \textbf{Clinical Examples:} Hernia repair, breast biopsy or reconstruction, thyroidectomy, splenectomy (elective), and craniotomy.
  \item \textbf{SSI Risk (without antibiotics):} $<2.0\%$.
  \item \textbf{Antibiotic Prophylaxis:} Generally not indicated unless a foreign body (prosthetic implant, mesh, joint replacement) is introduced, or in high-risk patients.
\end{itemize}

\section*{Class II: Clean-Contaminated Wounds}
\textbf{Definition:} An operative wound in which the respiratory, alimentary, genital, or urinary tracts are entered under controlled conditions and without unusual contamination.
\begin{itemize}
  \item \textbf{Criteria:}
  \begin{itemize}
    \item Controlled entry into a mucosal-lined organ system.
    \item Minor break in sterile technique.
    \item No evidence of active infection or gross contamination.
  \end{itemize}
  \item \textbf{Clinical Examples:} Elective cholecystectomy, uncomplicated appendectomy, hysterectomy, tonsillectomy, and cystoscopy.
  \item \textbf{SSI Risk (without antibiotics):} $3.0\% - 11.0\%$.
  \item \textbf{Antibiotic Prophylaxis:} Single dose of preoperative antibiotics (e.g., Cefazolin) administered within 60 minutes prior to surgical incision.
\end{itemize}

\section*{Class III: Contaminated Wounds}
\textbf{Definition:} Open, fresh, accidental wounds, or operations with major breaks in sterile technique or gross spillage from the gastrointestinal tract, or incisions in which acute, non-purulent inflammation is encountered.
\begin{itemize}
  \item \textbf{Criteria:}
  \begin{itemize}
    \item Gross spillage of gastrointestinal contents.
    \item Major break in aseptic technique (e.g., open cardiac massage).
    \item Traumatic wounds less than 4 hours old.
  \end{itemize}
  \item \textbf{Clinical Examples:} Inflamed appendix with localized spill, open fractures, and diverticulitis with localized purulent inflammation.
  \item \textbf{SSI Risk (without antibiotics):} $10.0\% - 17.0\%$.
  \item \textbf{Antibiotic Prophylaxis:} Preoperative antibiotics required; may be continued for 24 hours postoperatively. Delayed primary closure may be considered.
\end{itemize}

\section*{Class IV: Dirty or Infected Wounds}
\textbf{Definition:} Old traumatic wounds with retained devitalized tissue or those that involve existing clinical infection or perforated viscera.
\begin{itemize}
  \item \textbf{Criteria:}
  \begin{itemize}
    \item Organism causing postoperative infection was present in the operative field before the surgery.
    \item Perforated viscera (e.g., bowel, appendix) with gross fecal contamination.
    \item Traumatic wounds more than 4 hours old with necrotic tissue.
  \end{itemize}
  \item \textbf{Clinical Examples:} Perforated appendicitis with diffuse purulent peritonitis, abscess debridement, and gangrenous bowel resection.
  \item \textbf{SSI Risk (without antibiotics):} $>27.0\%$.
  \item \textbf{Therapeutic Management:} Represents an active infection. The administration of antibiotics is considered therapeutic (course of 3-7 days or until clinical resolution) rather than prophylactic. Primary skin closure is often avoided (wound left open to heal by secondary intention or delayed primary closure).
\end{itemize}

\clearpage
